\begin{table}[htbp]
\centering
\caption{Standardized Effect Sizes}
\label{tab:sde}
\small
\begin{tabular}{l cccccc}
\hline\hline
Outcome & $\hat{\beta}$ & SE & SD($Y$) & SDE & SE(SDE) & Classification \\
\hline
Total CO2e & -0.1457 & 0.0397 & 1.7275 & -0.0843 & 0.0230 & Mod.\ neg. \\
CO2 & -0.1636 & 0.0826 & 2.8537 & -0.0573 & 0.0289 & Mod.\ neg. \\
CH4 (CO2e) & 0.0030 & 0.0883 & 3.3365 & 0.0009 & 0.0265 & Null \\
N2O (CO2e) & -0.1200 & 0.0795 & 2.4318 & -0.0494 & 0.0327 & Small neg. \\
Energy-int.\ sectors & -0.2092 & 0.0738 & 1.8162 & -0.1152 & 0.0406 & Mod.\ neg. \\
Non-energy sectors & 0.1266 & 0.1074 & 1.2595 & 0.1005 & 0.0853 & Mod.\ pos. \\
\hline\hline
\end{tabular}
\begin{minipage}{0.95\textwidth}
\footnotesize
\textit{Notes:} Standardized effect sizes computed as SDE $= \hat{\beta} / \text{SD}(Y)$ for binary treatment.
Research question: Does mandatory carbon pricing reduce facility-level industrial greenhouse gas emissions?
Data: ECCC GHGRP, 2004--2023 ($N = 11,038$ facility-years). Method: TWFE DiD with facility and year fixed effects.
Treatment: federal carbon pricing backstop imposed on ON, SK, MB, NB (April 2019).
Classification labels refer to the magnitude of the standardized point estimate, not to statistical significance. ``Null'' denotes a near-zero effect size ($|$SDE$| < 0.005$), not a failure to reject a null hypothesis.
SDE $< -0.15$: large negative; $[-0.15, -0.05)$: moderate negative; $[-0.05, -0.005)$: small negative;
$[-0.005, 0.005]$: null; $(0.005, 0.05]$: small positive; $(0.05, 0.15]$: moderate positive; $> 0.15$: large positive.
\end{minipage}
\end{table}
